%!TEX root = ../scivis_lbaakman_bvanloon.tex

\chapter{Color Mapping} % (fold)
\label{cha:color_mapping}
Color mapping is widely used for the visualization of scalar data. This is done by associating a color with each scalar value that is visualized. In this chapter we will discuss the colormaps included in the application and how these colormaps are supported. In \cref{s:colormapping:method} the general methods used for colormapping is first explained as well as the parametrization and application of colormaps. Next \Cref{s:colormaps:differentmaps} highlights the important features of colormaps and discusses the colormaps that are implemented in our application. For each colormap that is provided the advantages and disadvantages are discussed, as well as a small motivation why it is included in the application. Finally, \cref{s:colormapping:results} shows the results of applying colormaps to the simulation and discusses these results.


\section{Method}
\label{s:colormapping:method}
%!TEX root = ../../scivis_lbaakman_bvanloon.tex
\section{Method}
\label{s:streamsurfaces:method}
This section explains how one can define a stream surface given a seed curve \seedCurve, with $n$ vertices, $\seedCurveVertex{0}, \cdots ,\seedCurveVertex{n}$ and $T$ two-dimensional vector fields, $\vectorFieldDiscTime{}{0}, \cdots, \vectorFieldDiscTime{}{T}$,  that represent the different states of some vector field as a function of time.


In \cref{s:streamsurface:method:time} we explain how we compute a three dimensional streamline in a number of consecutive two dimensional datasets. \Cref{s:streamsurface:method:seedCurve} explain how we compute the seed points given the seed curve \seedCurve. \Cref{s:streamsurface:method:surface} discusses how we define a stream surface from a set of streamlines. 

\subsection{Time as the Third Dimension}
\label{s:streamsurface:method:time}
We start the streamline at seed point \seedPoint{0} in the two-dimensional vector field \vectorFieldDiscTime{}{T}, \ie the oldest vector field that we have. We then compute the next vertex of the two-dimensional streamline in the vector field \vectorFieldDiscTime{}{T} with seed point \seedPoint{0}. We refer to this two-dimensional position as \seedPoint{1}. This position is used as the seed point of a two-dimensional streamline of one edge in the vector field \vectorFieldDiscTime{}{T - 1}. The other vertex of that edge is referred to as \seedPoint{2}, and is used as the seed point in \vectorFieldDiscTime{}{T - 2}. This process is repeated until we have reached \vectorFieldDiscTime{}{0}. 

The points $\seedPoint{0}, \cdots, \seedPoint{T}$ define the three dimensional stream line. The \textit{z}-coordinate of these points is computed in such a way that the points are spread uniformly in \textit{z} over some predefined range for \textit{z}.


The streamline is thus computed according to the process described in \cref{cha:streamlines}, with one change: there is no lower bound for the magnitude, \ie $\magnitudeMin = \infty$.

\subsection{From a Seed Curve to Seed Points}
\label{s:streamsurface:method:seedCurve}
The seed curve is a polyline, each of the line segments of this polyline is split into \resolution segments. The higher \resolution is the better the final stream surface is. 

\subsection{Building the Surface}
\label{s:streamsurface:method:surface}
The final step of defining a stream surface is the determination of a surface from a set of streamlines, \streamLineIdx{0}, \streamLineIdx{N}. The line segments of the stream lines are have all approximately the same length, as the distance between \seedPoint{i} and \seedPoint{i + 1} is equal for every segment in every stream line and the difference in \textit{z}-coordinates of consecutive stream line vertices is also equal. We then define the stream surface by connecting vertex \seedPoint{i} of stream line \streamLineIdx{j} with \seedPoint{i} of stream line \streamLineIdx{j + 1}, for $j < N$. If stream line \streamLineIdx{j} has a vertex \seedPoint{i}, but \streamLineIdx{j + 1} is shorter and does not have a \seedPoint{i}, we connect \seedPoint{i} of \streamLineIdx{j} to the last vertex of \streamLineIdx{j + 1}.


The resulting connections define a collection of quads that could define a surface. However the thus defined surface would happily move though some obstacle that is encountered by the flow of the vector field. To avoid this we do not connect two neighboring vertices of streamlines if the distance between them is greater than \divergenceCriterion. 


\section{Provided Color Maps}
\label{s:colormaps:differentmaps}
%!TEX root = ../scivis_lbaakman_bvanloon.tex
This section discusses the colormaps provided in the application. Each colormap is accompanied with a short discussion giving the main goal of the colormap and its pros and cons.

As discussed in the previous section, colormaps provide a way to map scalar values to colors in order for the user to extract information. The quality of a colormap can be judged by how well and intuitive the colormap allows inverse mapping. Which colormap is best used depends on the context and goal of the visualization. The zebra-colormap is for example well suited to visualize value variations (i.e. the first derivative of a scalar value), but would be less useful would the aim be to highlight maxima. Historical reasons can also influence the choice of colormaps. X-rays are for example standard displayed using a gray-scale colormap. Better alternatives might exist, but are not used due to status quo in the medical field.

In our application we have chosen to implement a variety of colormaps, this way the application offers a suitable colormap for almost any form of visualization, ranging from the (in)famous rainbow colormap to luminance based, two-hue and diverging colormaps.

\begin{figure}[b]
	\begin{subfigure}{1\textwidth}
		\begin{minipage}[l]{0.05\textwidth}
	    	\caption{} \label{fig:colormaps:legends:rainbow}
	  	\end{minipage}
	  \hfill
	  	\begin{minipage}[l]{0.95\textwidth}
	    	\includegraphics[width=\textwidth,height=15px,keepaspectratio=false,frame]{colormapping/img/colormap_legends/rainbowcolormap.png}
	  	\end{minipage}
	\end{subfigure}

  	\begin{subfigure}{1\textwidth}
  		\begin{minipage}[l]{0.05\textwidth}
    		\caption{} \label{fig:colormaps:legends::zebra}
  		\end{minipage}
  	\hfill
  		\begin{minipage}[l]{0.95\textwidth}
    		\includegraphics[width=\textwidth,height=15px,keepaspectratio=false,frame]{colormapping/img/colormap_legends/zebracolormap.png}
  		\end{minipage}
  	\end{subfigure}

  	\begin{subfigure}{1\textwidth}
  		\begin{minipage}[l]{0.05\textwidth}
    		\caption{} \label{fig:colormaps:legends::grayscale}
  		\end{minipage}
  	\hfill
  		\begin{minipage}[l]{0.95\textwidth}
    		\includegraphics[width=\textwidth,height=15px,keepaspectratio=false,frame]{colormapping/img/colormap_legends/grayscalecolormap.png}
  		\end{minipage}
  	\end{subfigure}
  	\begin{subfigure}{1\textwidth}
  		\begin{minipage}[l]{0.05\textwidth}
    		\caption{} \label{fig:colormaps:legends::heat}
  		\end{minipage}
  	\hfill
  		\begin{minipage}[l]{0.95\textwidth}
    		\includegraphics[width=\textwidth,height=15px,keepaspectratio=false,frame]{colormapping/img/colormap_legends/heatcolormap.png}
  		\end{minipage}
  	\end{subfigure}
  	\begin{subfigure}{1\textwidth}
  		\begin{minipage}[l]{0.05\textwidth}
    		\caption{} \label{fig:colormaps:legends::cold}
  		\end{minipage}
  	\hfill
  		\begin{minipage}[l]{0.95\textwidth}
    		\includegraphics[width=\textwidth,height=15px,keepaspectratio=false,frame]{colormapping/img/colormap_legends/coldcolormap.png}
  		\end{minipage}
  	\end{subfigure}
  	\begin{subfigure}{1\textwidth}
  		\begin{minipage}[l]{0.05\textwidth}
    		\caption{} \label{fig:colormaps:legends::hue}
  		\end{minipage}
  	\hfill
  		\begin{minipage}[l]{0.95\textwidth}
    		\includegraphics[width=\textwidth,height=15px,keepaspectratio=false,frame]{colormapping/img/colormap_legends/huecolormap.png}
  		\end{minipage}
  	\end{subfigure}
  	\begin{subfigure}{1\textwidth}
  		\begin{minipage}[l]{0.05\textwidth}
    		\caption{} \label{fig:colormaps:legends::twocolors}
  		\end{minipage}
  	\hfill
  		\begin{minipage}[l]{0.95\textwidth}
    		\includegraphics[width=\textwidth,height=15px,keepaspectratio=false,frame]{colormapping/img/colormap_legends/twocolorscolormap.png}
  		\end{minipage}
  	\end{subfigure}
  	\begin{subfigure}{1\textwidth}
  		\begin{minipage}[l]{0.05\textwidth}
    		\caption{} \label{fig:colormaps:legends::diverging}
  		\end{minipage}
  	\hfill
  		\begin{minipage}[l]{0.95\textwidth}
    		\includegraphics[width=\textwidth,height=15px,keepaspectratio=false,frame]{colormapping/img/colormap_legends/divergingcolormap.png}
  		\end{minipage}
  	\end{subfigure}
\caption{Colormap legends provided in the application: 
(\subref{fig:colormaps:legends:rainbow}) rainbow, 
\subref{fig:colormaps:legends::zebra} zebra,
\subref{fig:colormaps:legends::grayscale} gray-scale,
\subref{fig:colormaps:legends::heat} heat,
\subref{fig:colormaps:legends::cold} cold,
\subref{fig:colormaps:legends::hue} user picked hue,
\subref{fig:colormaps:legends::twocolors} linear two-color, and
\subref{fig:colormaps:legends::diverging} blue-red diverging. Each colormap is shown using full saturation and 256 colors (except for the zebra map which has 50 color regions). 
 }\label{fig:colormaps:legends}
\end{figure}


\subsection{Rainbow Colormap} % (fold)
\label{ssub:rainbow_colormap}
The rainbow colormap (given in \cref{fig:colormaps:legends:rainbow}) is one of the most used colormaps in the scientific community and is therefore included in our application. The colormap is often seen as visual appealing due to the varying colors and can give an intuitively mapping of high and low values to warm (red) and cold (blue) colors. However, although often used, the rainbow colormap has many flaws and is heavily critiqued\cite{borland2007rainbow}\cite{divergingMoreland2009}. A few of the most serious flaws are discussed below.
\begin{itemize}
  \item The rainbow colormap pulls the focus of users towards areas with high values (and thus warm, red colors). This can be an advantage when the area of interest is the maximum values in the visualization but might distract otherwise.
  \item The colors in the rainbow colormap have the same luminance (except for the tips of the colormap) but might be perceived to have different luminance by the user. As a result the ordering of the luminance in the colormap might appear to be non-monotonic which can confuse the inverse mapping process.
  \item There is no decisive ordering of the colors in the rainbow colormap. The colormap orders the colors from blue to green to yellow to red, but this is not a natural ordering perceived by humans.
  \item The colormap does not behave well when the saturation is lowered or when translated to black and white (when for example a copy is made). Doing so results in a broken, non-monotonic colormap disallowing proper inverse mapping.
  \item The colormap is not colorblind friendly.
\end{itemize}
% subsubsection flaws (end)


\subsection{Luminance Colormaps} % (fold)
\label{sub:luminance_colormaps}
The next set of colormaps use luminance to give a natural ordering of the colors (and thus values). These colormaps are often effective since human perception is sensitive for change in luminance\cite{divergingMoreland2009} which makes inverse mapping easier. There are three major disadvantages to luminance based colormaps. The colormaps are not well-suited for use in combination with 3D shading, since shading can cause change in the perceived luminance of a color or vice verse change in luminance might give the appearing of non-existing shading. Secondly luminance based colormaps draw attention to values with high luminance, details in low-luminance areas are often hard to see and will therefore be ignored by the user. A third flaw is that pixels with the same luminance might look different depending on the luminance of the surrounding pixels. 
% subsection luminance_colormaps (end)

\subsubsection{Gray-Scale colormap} % (fold)
\label{ssub:gray_scale_colormap}
The gray-scale colormap is the simplest of the luminance based colormaps and does not use any hue. The advantage of this colormap is that it will often be independent of the used medium (screens, paper, black-and-white print) and does not use the does not rely on distinguishing between hues for inverse mapping. This makes the colormap suitable for many application and users. Besides the general disadvantages associated with luminance based colormaps it has as flaw that it is not well suited for non compact datasets (e.g. datasets containing holes). In these cases the background might be similar to objects in the dataset making them blend into the background.


\subsubsection{Single Hue Colormap} % (fold)
\label{ssub:single_hue_colormap}
The single hue colormap is an extension on the gray-scale colormap to which a single user picked hue is added. This can make the colormap more visual appealing and removes the probability of the visualization blending in with the background.

\subsubsection{Heat Colormap} % (fold)
\label{sub:heat__colormap}
The heatmap is a luminance based colormap that uses both luminance and a ranges of hues for the ordering and thus provides more clues for the ordering0 compared to the gray-scale and single hue colormap. Although it uses less hues than the rainbow colormap it has as advantage that the used hues have a natural perceived ordering red-orange-yellow.

In some application the heat colormap uses white for the highest value data. Since our application has a white background, the white value is removed from the colormap making the color for the highest value data yellow. This makes the colormap suitable for visualizations as glyphs and streamlines, which might blend into the white background would the colormap contain white.


\subsubsection{Cold Colormap} % (fold)
\label{ssub:cold_colormap}
The cold colormap is similar to the heat colormap except it uses blue tinted colors instead of red tints. It is added since blue and red are nice contrasting colors making it possible to combine scalar visualization with glyphs by using the heatmap for one visualization and the cold colormap for the second visualization.

\subsection{Isoluminant Colormap} % (fold)
\label{sub:two_hue_colormap}
The isoluminant colormap is a cyan-to-mauve colormap based on linear interpolation based on \cite{divergingMoreland2009}. Because the colormap is isoluminant it is well suited to be used in combination with 3D shading since difference in luminance can only be caused by shading. Opposed to the green-to-red isoluminant colormap that is also often used, the cyan-to-mauve colormap is colorblind friendly, making it a more suitable candidate. Although well suited for use in 3D visualizations the colormap has two major flaws. The human perception is less sensitive to changes in hue than in changes in luminance, making the inverse mapping processes of the colormap harder. Furthermore, since the colormap offers less color variation than a lot of other colormaps making it more difficult to visualize a large range of values (and is the colormap often perceived as dull).
% subsection two_hue_colormap (end)

\subsection{Diverging Colormap} % (fold)
\label{sub:diverging_colormap}


+ Better perceive values diverging from the average
\subsection{Zebra Colormap} % (fold)
\label{sub:zebra_colormap}


\subsection{Results}
\label{s:colormapping:results}
%!TEX root = ../../scivis_lbaakman_bvanloon.tex
\section{Results}
\label{s:streamsurfaces:results}
This section presents and discusses several visualizations of the evolution of a two-dimensional vector field as a function of time with stream surfaces. All visualizations in this section will show the fluid velocity, and will be colored according to the fluid velocity magnitude with a color map that is not clamped. The older simulation states are shown near the bottom of the images, more recent states are shown near the top of the images.

\begin{figure}
	\centering
	\begin{subfigure}{0.7\textwidth}
		\centering
		\assusScreenshot[height=0.27\textheight, keepaspectratio=true]{./img/streamsurfaces/vertices}
		\caption{Vertices}
		\label{fig:streamsurfaces:differentmethods:vertices}
	\end{subfigure}
	\begin{subfigure}{0.7\textwidth}
		\centering
		\assusScreenshot[height=0.27\textheight, keepaspectratio=true]{./img/streamsurfaces/lines}
		\caption{Stream Lines}
		\label{fig:streamsurfaces:differentmethods:lines}
	\end{subfigure}	
	\begin{subfigure}{0.7\textwidth}
		\centering
		\assusScreenshot[height=0.27\textheight, keepaspectratio=true]{./img/streamsurfaces/surface}
		\caption{Stream Surface}
		\label{fig:streamsurfaces:differentmethods:surface}
	\end{subfigure}		
	\caption{Visualization of the change of a two-dimensional vector field, \velocity, as function of time with \subref{fig:streamsurfaces:differentmethods:vertices} the vertices of the stream lines, \subref{fig:streamsurfaces:differentmethods:lines} the stream lines and \subref{fig:streamsurfaces:differentmethods:surface} the stream surface. %
	\resolution = 10, \numStates = 200, \divergenceCriterion = 200.0.}
	\label{fig:streamsurfaces:differentmethods}
\end{figure}

We support three different ways of visualizing a stream surface. Firstly one can use the vertices of the streamlines as shown in \cref{fig:streamsurfaces:differentmethods:vertices}. Alternatively one can visualize the streamlines as in \cref{fig:streamsurfaces:differentmethods:lines}, or the surface defined by these lines as in \cref{fig:streamsurfaces:differentmethods:surface}. 

\Cref{fig:streamsurfaces:differentmethods:surface} clearly illustrates one issues of our stream surfaces, firstly the edges of the computational domain are not reached gracefully. It is possible that this is caused by the fact that our streamlines end at their last vertex before the end of the computational domain, instead of on the edge of the domain.

It should be noted that the resolution of seed lines used for the visualizations in \cref{fig:streamsurfaces:differentmethods} was chosen quite low, to emphasize the difference between the lines and the vertices.

The straight streamlines at the bottom of these figures, \ie in the oldest simulation states, are caused by lack of user input at the beginning. In that case, all vectors are \vec{0}, consequently $\seedPoint{i} = \seedPoint{i + 1}$, which results in a line segment whose vertices only differ in their \textit{z}-values.

\Cref{fig:streamsurfaces:differentResolution} presents three visualizations of the same data set with three different resolutions. This figures clearly illustrates the importance of choosing a resolution that is sensible for the data, in this case a resolution of five is clearly too low, \resolution = 50, seems sufficient. If resolution is set to 100 the stream lines are so close together that they nearly form a surface. 

\begin{figure}
	\centering
	\begin{subfigure}[b]{0.3\textwidth}
		\centering
		\threeDScreenshot{./img/streamsurfaces/r5}
		\caption{\resolution = 5}
		\label{fig:streamsurfaces:differentResolution:r5}
	\end{subfigure}
	\begin{subfigure}[b]{0.3\textwidth}
		\centering
		\threeDScreenshot{./img/streamsurfaces/r50}
		\caption{\resolution = 50}
		\label{fig:streamsurfaces:differentResolution:r50}
	\end{subfigure} 
	\begin{subfigure}[b]{0.3\textwidth}
		\centering
		\threeDScreenshot{./img/streamsurfaces/r100}
		\caption{\resolution = 100}
		\label{fig:streamsurfaces:differentResolution:r100}
	\end{subfigure} 
	\caption{Visualization of the change of a two-dimensional vector field, \velocity, as function of time with different resolutions. %
	\numStates = 200}
	\label{fig:streamsurfaces:differentResolution}
\end{figure}

The influence of the \divergenceCriterion is illustrated in \cref{fig:streamsurfaces:differentDivergenceCriterion}. To more clearly show the divergence both the stream lines and the stream surfaces are shown. Since the length of the stream lines are not influenced by \divergenceCriterion they are always visible independent of the divergence of the vector field.

The influence of the \divergenceCriterion can be seen in the size of the triangle that is placed between the flow that moves to the left and the flow to the right. The higher the divergence criterion the longer the two flows are connected. If \divergenceCriterion = 3 the stream splits quite early, whereas for \divergenceCriterion = 200 the streams stay together until \rfrac{1}{3} of the y-axis. The influence of the divergence criterion is also illustrated in the upper left corner of these images. 

\begin{figure}
	\centering
	\begin{subfigure}[b]{0.3\textwidth}
		\centering
		\threeDScreenshot{./img/streamsurfaces/d3}
		\caption{\resolution = 3}
		\label{fig:streamsurfaces:differentDivergenceCriterion:r3}
	\end{subfigure}
	\begin{subfigure}[b]{0.3\textwidth}
		\centering
		\threeDScreenshot{./img/streamsurfaces/d30}
		\caption{\resolution = 30}
		\label{fig:streamsurfaces:differentDivergenceCriterion:r30}
	\end{subfigure} 
	\begin{subfigure}[b]{0.3\textwidth}
		\centering
		\threeDScreenshot{./img/streamsurfaces/d200}
		\caption{\resolution = 200}
		\label{fig:streamsurfaces:differentDivergenceCriterion:r200}
	\end{subfigure} 
	\caption{Visualization of the change of a two-dimensional vector field, \velocity, as function of time with different divergence criteria. %
	\resolution = 100, \numStates = 200.}
	\label{fig:streamsurfaces:differentDivergenceCriterion}
\end{figure}

\Cref{fig:streamsurfaces:differentVectorFields} shows the stream lines of two different vector fields, namely the fluid velocity field and the force field. \Cref{fig:streamsurfaces:differentVectorFields:force} clearly illustrates that the added forces decays quite quickly. This visualization also shows that the stream lines are straight lines as long as the vector fields don't change as a function of time at the position of the seed point.

\begin{figure}
	\centering
	\begin{subfigure}[b]{0.3\textwidth}
		\centering
		\threeDScreenshot{./img/streamsurfaces/fluidVelocity}
		\caption{\velocity}
		\label{fig:streamsurfaces:differentVectorFields:velocity}
	\end{subfigure}
	\begin{subfigure}[b]{0.3\textwidth}
		\centering
		\threeDScreenshot{./img/streamsurfaces/force}
		\caption{\force}
		\label{fig:streamsurfaces:differentVectorFields:force}
	\end{subfigure}  
	\caption{Visualization of the change of \subref{fig:streamsurfaces:differentVectorFields:velocity} fluid velocity field and \subref{fig:streamsurfaces:differentVectorFields:force} the force field as a function of time. 
	%
	\resolution = 100, \numStates = 200}
	\label{fig:streamsurfaces:differentVectorFields}
\end{figure}

% chapter color_mapping (end)